% Generated draft: requires mathematical and source review.
\documentclass[11pt]{article}
\usepackage[margin=25mm]{geometry}
\usepackage{fontspec}
\setmainfont{Latin Modern Roman}
\usepackage{amsmath,amssymb,mathtools,unicode-math}
\setmathfont{Latin Modern Math}
\usepackage{enumitem,needspace,xurl,hyperref}
\setlength{\parindent}{0pt}
\setlength{\parskip}{5pt}
\setlength{\emergencystretch}{3em}
\newcommand{\N}{\mathbb{N}}
\newcommand{\Z}{\mathbb{Z}}
\newcommand{\Q}{\mathbb{Q}}
\newcommand{\R}{\mathbb{R}}
\newcommand{\C}{\mathbb{C}}
\newcommand{\F}{\mathbb{F}}
\newcommand{\A}{\mathbb{A}}
\newcommand{\PP}{\mathbb{P}}
\newcommand{\E}{\mathbb{E}}
\newcommand{\eps}{\varepsilon}
\newcommand{\dd}{\,\mathrm{d}}
\begin{document}
\section*{UnsolvedMath ID: 20002901}
\textbf{Source problem code:} catalogue entry
\subsection{Definitions and mathematical statement}
The source asks to generalize \(k\)-Schur functions beyond root-system type A. It does not specify which properties must be retained: a symmetric-function realization, affine-Schubert basis, Hopf structure, positivity, or a Macdonald-type parameter deformation. These are distinct versions of the request.

The supplied aid reports an established geometric specialization at deformation parameter \(t=1\) for every crystallographic root system: the affine-Grassmannian Schubert basis, equivalently its triangular Peterson/Fomin--Stanley realization. For a simply connected semisimple complex group \(G\), its homology \(H_*(\operatorname{Gr}_G;\Z)\) has a Schubert basis indexed by affine Grassmannian elements of the affine Weyl group. In type \(A_r\), \(G=\mathrm{SL}_{r+1}\), this recovers the \(k=r\) Schur basis at \(t=1\) in
\[
 \Lambda_{(r)}=\Z[h_1,\ldots,h_r],\qquad\deg h_i=i.
\]
Use half the topological homological grading throughout. This is not a supplied solution for noncrystallographic root systems or every parameter-dependent \(k\)-Schur theory. The aid reports exact symmetric-function models in types C and B, while a uniform Macdonald-type deformation remains incomplete in that attempt.

\textbf{Root-degree obstruction reported in the aid.} Let the root system have rank \(r\) and exponents \(m_1,\ldots,m_r\). Rationally, the affine-Grassmannian-compatible homology algebra has polynomial generator degrees \(m_i\), hence
\[
 \operatorname{Hilb}(H_*(\operatorname{Gr}_G;\Q);u)
           =\prod_{i=1}^r(1-u^{m_i})^{-1}.
\]
Equivalently, its rational augmentation indecomposables have precisely those degrees, with multiplicities. If a proposed generalization used the unchanged graded type A algebra \(\Q[h_1,\ldots,h_r]\) with degrees \(1,\ldots,r\), equality of graded Hilbert series would require
\[
 \{m_1,\ldots,m_r\}=\{1,\ldots,r\}.
\]
For finite crystallographic root systems this forces type \(A_r\), up to the usual low-rank isomorphisms. Reducible systems have one exponent 1 per component and must retain their multiplicities. The obstruction is graded and root-sensitive; it is not a claim that abstract ungraded polynomial algebras in \(r\) variables are nonisomorphic.

The aid's integral type \(C_2\) witness uses one-row Schur \(P\)-functions. Normalize them by
\[
 \prod_{a\ge1}\frac{1+x_a u}{1-x_a u}
       =1+2\sum_{j\ge1}P_j u^j.
\]
In the stated type C model the ring is
\[
 \Z[P_1,P_3],\qquad\deg P_1=1,\quad\deg P_3=3,
\]
with
\[
 P_2=P_1^2,\qquad P_4=2P_1P_3-P_1^4.
\]
Thus \(P_2\) cannot serve as an independent degree-two generator, and the correct exponent profile is \((1,3)\). These relations provide an integral illustration, not a uniform construction for all other types or a Macdonald deformation.

The remaining formulation task is to specify which representation-theoretic and combinatorial properties define the desired non-A functions and how a parameter \(k\) and deformation variables should vary. The geometric \(t=1\) basis and exponent constraint answer parts of that request without selecting a unique general theory.
\subsection{Short English statement}
The aid reports affine-Grassmannian Schubert homology as the crystallographic \(t=1\) generalization. Its generator degrees must follow root exponents, so the unchanged graded type A algebra works only in type A. A type \(C_2\) Schur-\(P\) example shows the required degree change, while a uniform parameter deformation remains outside the supplied result.
\subsection{Sources}
\begin{itemize}
\item Local UnsolvedMath registry, permanent ID 20002901, catalogue entry; canonical \texttt{aim-special-functions-notes.json}, notes[5].
\item AIM \emph{Generalized Kostka polynomials}, Problem 6, \url{https://aimath.org/WWN/kostka/prob2.pdf}. No individual proposer or authoritative bibliography is supplied.
\item The aid names Peterson/Fomin--Stanley realizations and type B/C symmetric-function models without full titles, authors for those models, dates or theorem references.
\item Machine-generated research attempt 1, partial-result/partially-solved labels, reported rigor proved and low novelty confidence. Its exponent-profile obstruction and integral type \(C_2\) illustration are not peer reviewed and require independent verification.
\end{itemize}
\end{document}
